\section{Exact algebra and complete proof map}
\label{app:proofs}

This appendix makes every finite calculation used above explicit.  The
matrix convention and multiplicity input are source-owned
\citep{knauf1998spin,knauf1999erratum,knauf1993phases}; all deductions from
them are reproduced here.

\subsection{Recurrence and exact products}

For $M_u=\left(\begin{smallmatrix}a&b\\c&d\end{smallmatrix}\right)$,
right multiplication by $L$ preserves the first column, while right
multiplication by $R$ replaces it by the sum of the two columns.  The
conjugacy $M_{\bar u}=JM_uJ$ identifies the second-column sum with
$h(\bar u)$.  This proves
$h(u0)=h(u)$ and $h(u1)=h(u)+h(\bar u)$ without transposition or reversed word
order.

The exact products needed by all witnesses are
\begin{table}[H]
\centering
\caption{Complete finite witness ledger.}
\label{tab:exact-products}
\small
\begin{tabular}{@{}c c c@{}}
\toprule
$w$ & $M_w$ & $h(w)$ \\
\midrule
$\eps$ & $\left(\begin{smallmatrix}1&0\\0&1\end{smallmatrix}\right)$ & $1$ \\
$0$ & $\left(\begin{smallmatrix}1&1\\0&1\end{smallmatrix}\right)$ & $1$ \\
$1$ & $\left(\begin{smallmatrix}1&0\\1&1\end{smallmatrix}\right)$ & $2$ \\
$01$ & $\left(\begin{smallmatrix}2&1\\1&1\end{smallmatrix}\right)$ & $3$ \\
$10$ & $\left(\begin{smallmatrix}1&1\\1&2\end{smallmatrix}\right)$ & $2$ \\
$11$ & $\left(\begin{smallmatrix}1&0\\2&1\end{smallmatrix}\right)$ & $3$ \\
$001$ & $\left(\begin{smallmatrix}3&2\\1&1\end{smallmatrix}\right)$ & $4$ \\
$010$ & $\left(\begin{smallmatrix}2&3\\1&2\end{smallmatrix}\right)$ & $3$ \\
\bottomrule
\end{tabular}
\end{table}

\subsection{Stable quotient and rooted action}

Let $\sim_0$ be generated by $w\sim_0w0$.  Repeated use of the recurrence
shows that $h$ is constant on each equivalence class.  Since
$\eps\sim_00$, a descended right append-one action would imply
$[1]=[01]$.  The table gives $h(1)=2$ and $h(01)=3$, contradicting class
invariance.  The contradiction uses the same invariant that certifies the
source quotient; it does not assume a topology or an undeclared shift map.

\subsection{Clock and phase descent}

The words $01$ and $10$ differ by a cyclic rotation, but their rooted labels
are $3$ and $2$.  Hence no function on necklaces can equal $h$ on every
representative.  The power witness is independent: $11=1^2$ as a word, yet
$h(11)=3\ne4=h(1)^2$.  Because $h>0$, these identities are necessary for
$\log h$ to be representative-independent and additive under powers.

For the literal phase, $001$ and $010$ are rotations and have labels $4$ and
$3$.  Their Liouville values are $+1$ and $-1$.  Also
$\lambda(h(11))=\lambda(3)=-1$, whereas
$\lambda(h(1))^2=\lambda(2)^2=+1$.  Finally, a one-letter character agreeing
on $0$ and $1$ must have generator values $1$ and $-1$, so it takes the
value $+1$ on $11$, contradicting the exact value $-1$.

\subsection{Inventory determinant and repair trilemma}

The source multiplicity theorem gives one eigenvalue $n^{-s}$ for each of
the $\varphi(n)$ stable states with label $n$.  Absolute summability on
$\RePart s>2$ proves trace class, justifies grouping the trace, and yields the
product in \Cref{thm:inventory-determinant}.  For $|u|<1$, the norm-convergent
logarithm gives the trace powers.  This chain uses no cyclic-word
identification.  In particular, the $n=1$ inventory state forces the factor
$1-u$.

The declared source-preserving repair list is now exhausted.  Quotienting rooted words by
rotation contradicts the clock witness.  Keeping the literal scalar phase
contradicts the phase witnesses.  Keeping the stable quotient with the full
right binary action contradicts append-one non-descent.  Using $Q_s$ changes
the marker to inventory powers.  Trace/eigenvalue clocks and enlarged matrix
states change the source object or clock.  Thus at least one frozen
coordinate must be surrendered; the statement does not range over
undeclared extensions.

The positive controls confirm that the proof does not prohibit cyclic or
temporal matrix data in general.  Direct calculation gives
$\Tr(LR)=\Tr(RL)=3$, and the spectral radius of a fixed positive matrix obeys
$\rho(M^r)=\rho(M)^r$.  Both controls pass only after replacing the rooted
first-column-sum clock.
